\chapter{Menangkap Sinyal: Estimasi, Filter, dan Radar}
\label{ch:sinyal}

Pemrosesan sinyal masuk ke perjalanan saya lewat dua mata kuliah. Kuliah optimum and adaptive
signal processing systems, dengan latihan kelompok berisi enam tugas, membahas teori estimasi,
filter optimal, dan filter adaptif. Pada musim dingin 2025/26, latihan radar system engineering
mengajak kami membangun radar di dalam notebook: persamaan radar, korelasi, FFT, deteksi, CFAR,
radar FMCW, dan beamforming digital, sampai sebuah simulator radar sederhana.

Pemrosesan sinyal sudah ``belajar dari data'' jauh sebelum istilah machine learning populer. Filter
adaptif yang meredam gema di telepon memperbarui bobotnya dengan gradien stokastik, persis seperti
jaringan saraf. Banyak konsep di bab ini, dari batas Cramér--Rao sampai Kalman filter, sudah kita
temui dari sisi machine learning. Di sini kita melihatnya dari sisi insinyur sinyal.

\section{Sinyal acak}

Sinyal yang menarik hampir selalu bercampur noise, sehingga kita memperlakukannya sebagai proses
acak. Sebuah proses disebut stasioner dalam arti luas kalau rata-ratanya tetap dan fungsi
autokorelasinya, $r_x[k]=\E\big[x[n]x[n+k]\big]$, hanya bergantung pada jarak waktu $k$, bukan pada
waktu absolut. Transformasi Fourier dari autokorelasi adalah \istilah{power spectral density}, yang
menyatakan bagaimana daya sinyal tersebar di seluruh frekuensi. Noise putih punya autokorelasi yang
hanya tidak nol di $k=0$, sehingga spektrumnya datar: semua frekuensi mendapat daya yang sama,
seperti cahaya putih yang memuat semua warna.

\section{Estimasi klasik}

Kuliah teori estimasi dimulai dengan contoh paling sederhana: menaksir level tegangan tetap $A$
dari pengukuran yang berisik, $x[n]=A+w[n]$ untuk $n=0,\dots,N-1$, dengan $w[n]$ noise putih Gauss
bervarians $\sigma^2$. Rata-rata sampel tak bias, variansnya $\sigma^2/N$, dan batas Cramér--Rao
(\cref{ch:mlat}) untuk masalah ini juga $\sigma^2/N$. Jadi rata-rata sampel adalah estimator
\istilah{minimum variance unbiased}, dan tidak ada estimator tak bias lain yang lebih baik
\citep{kay1993}.

Untuk model yang lebih umum, beberapa jalan tersedia. Kalau model linear dalam parameter,
$\vx=\bm H\bm\theta+\bm w$, estimator \istilah{best linear unbiased} memberi varians terkecil di antara
semua estimator linear tak bias, dan untuk noise Gauss ia sama dengan MVU. Estimator kuadrat
terkecil, $\hat{\bm\theta}=(\bm H^\top\bm H)^{-1}\bm H^\top\vx$, tidak memerlukan asumsi statistik apa
pun, hanya meminimalkan selisih kuadrat antara data dan model. Maximum likelihood menghubungkan
semuanya: untuk noise Gauss, ia menghasilkan kuadrat terkecil, dan untuk sampel besar ia mendekati
batas Cramér--Rao.

\section{Estimasi Bayesian}

Kalau kita punya pengetahuan awal tentang parameter, misalnya bahwa sebuah tegangan hampir pasti di
sekitar lima volt, kita memperlakukan parameter sebagai variabel acak dengan prior. Estimator yang
meminimalkan galat kuadrat rata-rata, \istilah{minimum mean square error}, ternyata adalah rata-rata
posterior $\E[\theta\mid\vx]$. Estimator MAP mengambil puncak posterior. Untuk model linear Gauss,
keduanya sama, dan rata-rata posterior adalah kompromi berbobot antara prior dan data: kalau data
sedikit atau berisik, hasilnya dekat ke prior, kalau data banyak, hasilnya mengikuti data. Ini
adalah logika yang sama dengan Kalman gain di \cref{ch:prob}, dan memang Kalman filter adalah
estimator MMSE yang berjalan berurutan dalam waktu.

\section{Filter optimal}

Misalkan kita ingin mendapatkan sinyal $d[n]$ dari pengamatan $x[n]$ yang terganggu, dengan sebuah
filter linear berbobot $\bm w$ yang bekerja pada $M$ sampel terakhir, $y[n]=\bm w^\top\vx[n]$. Filter
Wiener memilih bobot yang meminimalkan galat kuadrat rata-rata $\E\big[(d[n]-y[n])^2\big]$. Menurunkan
terhadap $\bm w$ dan menyamakannya dengan nol memberi persamaan Wiener--Hopf,
\begin{equation}
\bm R\,\bm w=\bm p,\qquad \bm w_{\mathrm{opt}}=\bm R^{-1}\bm p,
\end{equation}
dengan $\bm R=\E[\vx\vx^\top]$ matriks autokorelasi input dan $\bm p=\E[\vx d]$ korelasi silang
antara input dan sinyal yang diinginkan. Hasil ini punya tafsiran geometris yang indah, yaitu
\istilah{orthogonality principle}: pada bobot optimal, galat tegak lurus terhadap semua data yang
dipakai. Tidak ada lagi informasi dalam data yang bisa dipakai untuk mengurangi galat.

Filter Wiener membutuhkan statistik $\bm R$ dan $\bm p$, yang dalam praktik harus diperkirakan dari
data. Menggantinya dengan rata-rata sampel memberi filter kuadrat terkecil, yang optimal untuk
data yang ada, bukan untuk ekspektasinya \citep{manolakis2005}.

\section{Filter adaptif}

Ketika statistik sinyal berubah dari waktu ke waktu, misalnya gema di ruang yang orangnya berpindah
tempat, filter harus terus menyesuaikan diri. Algoritme LMS mengambil satu langkah gradient descent
pada galat kuadrat sesaat di setiap sampel,
\begin{equation}
\bm w[n+1]=\bm w[n]+\mu\,e[n]\,\vx[n],\qquad e[n]=d[n]-\bm w[n]^\top\vx[n].
\end{equation}
Ini SGD dengan mini-batch berukuran satu, lima puluh tahun sebelum deep learning. Langkah $\mu$ harus
cukup kecil agar stabil, kira-kira $0<\mu<2/\lambda_{\max}$ dengan $\lambda_{\max}$ nilai eigen terbesar
$\bm R$, dan kecepatan konvergensinya ditentukan oleh sebaran nilai eigen $\bm R$. Inilah masalah
condition number dari \cref{ch:optimasi}, dalam pakaian yang berbeda. Algoritme RLS memperbarui
perkiraan $\bm R^{-1}$ secara rekursif sehingga setiap langkah mirip langkah Newton. Ia konvergen
jauh lebih cepat dan tidak terpengaruh sebaran nilai eigen, dengan biaya sekitar $M^2$ operasi per
sampel, dibanding $M$ untuk LMS.

Bayangkan headphone peredam bising. Mikrofon luar menangkap suara mesin pesawat, dan filter adaptif
belajar memprediksi seberapa banyak suara itu bocor ke telinga. Speaker lalu memutar prediksi
dengan tanda terbalik. Setiap kali pesawat berbelok atau penumpang menggeser kepala, galat sisa
berubah, dan bobot filter ikut bergeser beberapa ribu kali per detik.

\section{Radar}

Radar memancarkan gelombang elektromagnetik dan mendengarkan pantulannya. Daya yang kembali
mengikuti persamaan radar,
\begin{equation}
P_r=\frac{P_tG^2\lambda^2\sigma}{(4\pi)^3R^4},
\end{equation}
dengan $P_t$ daya pancar, $G$ penguatan antena, $\lambda$ panjang gelombang, $\sigma$ \istilah{radar
cross section} target, dan $R$ jaraknya. Pangkat empat pada jarak datang dari dua perjalanan: gelombang
melemah sebanding $R^2$ dalam perjalanan pergi dan sekali lagi dalam perjalanan pulang. Akibatnya,
untuk melipatgandakan jangkauan, daya pancar harus dinaikkan enam belas kali.

Untuk menemukan pantulan lemah di tengah noise, sinyal yang diterima dikorelasikan dengan sinyal yang
dipancarkan. Filter ini, \istilah{matched filter}, memaksimalkan rasio sinyal terhadap noise untuk
noise putih \citep{richards2014}. Puncak korelasi menunjukkan waktu tunda, dan waktu tunda menunjukkan
jarak, $R=c\tau/2$.

Setelah itu muncul pertanyaan deteksi: apakah sebuah puncak benar-benar target, atau hanya noise
yang kebetulan tinggi? Kriteria Neyman--Pearson menetapkan ambang yang memaksimalkan peluang deteksi
untuk peluang alarm palsu tertentu. Masalahnya, tingkat noise dan clutter berubah dari satu tempat
ke tempat lain. \istilah{Constant false alarm rate} (CFAR) menghitung ambang secara lokal. Untuk setiap
sel yang diuji, tingkat noise diperkirakan dari sel-sel latih di sekitarnya, dengan beberapa sel
penjaga di antaranya agar target sendiri tidak ikut terhitung. Varian cell-averaging merata-ratakan
sel latih, sedangkan varian order-statistic memakai nilai berperingkat tertentu dan lebih tahan
ketika ada beberapa target berdekatan \citep{rohling1983}.

\subsection{Radar FMCW}

Radar mobil dan radar di dalam ponsel umumnya memakai gelombang kontinu yang frekuensinya naik secara
linear, sebuah \istilah{chirp}. Pantulan dari target pada jarak $R$ tiba dengan tunda $2R/c$, sehingga
pada setiap saat frekuensinya sedikit lebih rendah daripada sinyal yang sedang dipancarkan. Mencampur
keduanya menghasilkan \istilah{beat frequency} yang sebanding dengan jarak,
\begin{equation}
f_b=\frac{2RS}{c},\qquad S=\frac{B}{T_c},
\end{equation}
dengan $B$ lebar pita chirp dan $T_c$ durasinya. Dengan $B=1$ GHz dan $T_c=50\,\mu$s, kemiringannya
$2\times10^{13}$ Hz per detik, dan target pada sepuluh meter menghasilkan beat sekitar
$1{,}33$ MHz. FFT dari sinyal beat langsung memberi profil jarak. Resolusi jaraknya $c/(2B)$,
lima belas sentimeter untuk lebar pita satu gigahertz.

Kecepatan diperoleh dari deretan chirp. Target yang bergerak menggeser fase pantulan sedikit dari satu
chirp ke chirp berikutnya, sebanding dengan kecepatannya. FFT kedua di sepanjang chirp mengubah
geseran fase itu menjadi kecepatan, dan hasil kedua FFT bersama-sama membentuk peta jarak dan
kecepatan. Dengan beberapa antena penerima, perbedaan fase antar antena menunjukkan arah datangnya
pantulan, $\Delta\phi=2\pi d\sin\theta/\lambda$. FFT ketiga, atau \istilah{digital beamforming}, mengubahnya
menjadi sudut. Latihan kami berakhir dengan simulator yang menggabungkan semua langkah ini: memancarkan
chirp ke beberapa target, menambahkan noise, lalu menjalankan seluruh rantai pemrosesan untuk
menemukannya kembali.

\section{Di mana machine learning masuk}

Batas antara pemrosesan sinyal dan machine learning semakin kabur. Peta jarak dan kecepatan radar
diperlakukan sebagai gambar dan dimasukkan ke jaringan konvolusi untuk mengenali gerakan tangan atau
membedakan pejalan kaki dari sepeda. Detektor CFAR klasik bisa diganti detektor yang dipelajari dari
data, dengan risiko bahwa jaminan alarm palsunya hilang. Algoritme iteratif, misalnya untuk
rekonstruksi sinyal, bisa ``dibuka'' (\istilah{unrolled}): setiap iterasi menjadi satu lapisan jaringan
yang parameternya dipelajari, sehingga struktur algoritme klasik tetap terjaga. Dan di arah sebaliknya, gagasan pemrosesan sinyal, seperti filter, frekuensi, dan
wavelet, memberi bahasa untuk memahami apa yang dilakukan jaringan saraf (\cref{ch:gdl}). Fondasinya
tetap sama: model, noise, dan batas atas apa yang bisa diestimasi dari data.

\sumberbab{Bab ini mengikuti kuliah dan latihan Optimum and Adaptive Signal Processing Systems
(variabel dan proses acak, estimasi klasik dan Bayesian, filter Wiener dan kuadrat terkecil, filter
adaptif) serta latihan Radar System Engineering (RCS dan persamaan radar, korelasi, FFT, deteksi,
CFAR, FMCW, beamforming digital, simulator). Bacaan yang disebut kuliah: \citet{kay1993} dan
\citet{manolakis2005}. Untuk radar, \citet{richards2014}.}
